\documentclass[conference]{IEEEtran}
\IEEEoverridecommandlockouts
\usepackage{cite}
\usepackage{amsmath,amssymb,amsfonts}
\usepackage{graphicx}
\graphicspath{{../figs/}}
\usepackage{textcomp}
\usepackage{xcolor}
\usepackage{booktabs}
\usepackage{url}
\renewcommand{\topfraction}{0.95}
\renewcommand{\bottomfraction}{0.95}
\renewcommand{\textfraction}{0.05}
\renewcommand{\floatpagefraction}{0.9}
\renewcommand{\dbltopfraction}{0.95}
\renewcommand{\dblfloatpagefraction}{0.9}
\def\BibTeX{{\rm B\kern-.05em{\sc i\kern-.025em b}\kern-.08em
    T\kern-.166em\lower.7ex\hbox{E}\kern-.125emX}}

\begin{document}

\title{Clause-Activation Signatures (CAS) for Attack-Family Attribution in IoMT}

\author{Rakesh Reddy Yakkati, Per-Arne Andersen, Lei Jiao, Ole-Christoffer Granmo, and Rishad Shafik
\thanks{Supported by the Research Council of Norway (SecureIoTM, grant 342167, TEKNOLOGIKONVER), with support from CAIR at the University of Agder, S\o rlandet Hospital Trust, and Newcastle University.}%
\thanks{R. R. Yakkati, L. Jiao, and O.-C. Granmo: Dept. of ICT, University of Agder, Grimstad, Norway; P.-A. Andersen: Dept. of Information Systems, University of Agder, Kristiansand, Norway; R. Shafik: Dept. of Microelectronic Systems, Newcastle University, UK. Contact: rakeshry@uia.no.}}

\maketitle

\begin{abstract}
Clause-Activation Signatures (CAS) convert a trained weighted Tsetlin
Machine (TM) into signed evidence lists for attack-family attribution.
Complementary-pairs stability selection (CPSS) identifies clauses
repeatedly selected across complementary half-samples; coefficient
signs separate evidence for and against each class. Weight-based pruning reduces the candidate pool before
signature extraction. On capped CICIoMT2024 WiFi/MQTT data, a three-way
fit/selection/report reanalysis across three trained seeds gives macro-F1
of $0.692\pm0.004$ for the base TM, $0.765\pm0.029$ for full-pool CAS,
and $0.787\pm0.023$ after pruning from 2,400 to 390 clauses. Pruning
reduces threshold sensitivity while lowering base macro-F1 by 0.004.
In the binary task, CAS raises macro-F1 from $0.739\pm0.006$ to
$0.824\pm0.002$ and achieves AUROC $0.929\pm0.021$ versus
$0.874\pm0.029$ for the TM class-sum margin. We derive a pruning margin condition and state
the assumptions and scope of stability-selection error control.
The tested synthetic signature unions are exactly identifiable, but
this does not establish real-traffic mixture detection. The analysis is retrospective;
independent evaluation remains necessary.
\end{abstract}

\begin{IEEEkeywords}
attack attribution, intrusion detection, interpretable machine learning,
IoMT, stability selection, Tsetlin Machine
\end{IEEEkeywords}

\section{Introduction and Related Work}
Attack-family attribution identifies the type of malicious traffic,
helping analysts prioritize alerts and choose appropriate responses in
the Internet of Medical Things (IoMT). Neural-network approaches address
this task through multiclass intrusion classification; for example,
Yin et al.~\cite{yin2017} used recurrent neural networks to distinguish
normal traffic and four attack categories on NSL-KDD. However, these
predictions do not directly expose human-readable decision rules.
Post-hoc methods such as SHAP and LIME explain feature contributions to
intrusion predictions~\cite{xaiids2025}, but a predicted label alone
still lacks explicit evidence explaining why one attack family is
favored over another.

The Tsetlin Machine (TM) learns conjunctive rules over Boolean
literals~\cite{granmo2018}, while weighted variants assign learned
weights to these clauses~\cite{phoulady2020}. Its rule-based structure
makes the conditions contributing to a class vote available for
inspection. Previous studies have used TM rules for interpretable
intrusion detection~\cite{abeyrathna2020}, including recent IoMT
applications that examine clause activations and class
votes~\cite{jaiswal2026iomt,jaiswal2026medsec}. These studies motivate
using the detector's own Boolean clauses as evidence for attack-family
attribution. A trained clause pool, however, is not automatically a compact class
signature. It can contain redundant rules, and a clause's importance
to the TM vote need not reflect its usefulness for distinguishing one
family from the others. A useful signature should expose both
supporting and opposing evidence, remain reusable across flows, and
keep each score traceable to the clauses that fired. Compactness and
repeatable selection must therefore be assessed alongside predictive
performance, rather than treated as sufficient evidence of explanation
quality.

We introduce Clause-Activation Signatures (CAS), which apply
complementary-pairs stability selection
(CPSS)~\cite{meinshausen2010,shah2013} to a frozen clause pool and retain
signed evidence for and against each class. Optional weight-based
pruning reduces the candidate pool before signature extraction, and
normalized evidence matching produces the attribution score. Our
contributions are this attribution layer, a pruning margin condition,
and an empirical comparison of full and pruned pools.

\section{Method: Clause-Activation Signatures}
\subsection{Booleanisation and Weighted TM}
\label{sec:bool}
Figure~\ref{fig:cpss} summarizes the pipeline. Training-fitted quantile
bins encode numeric features as cumulative thermometer bits; non-finite
values are replaced with the finite training maximum. TMU~\cite{tmu}
supplies literal negation internally. For $C$ classes and $K$ clauses
per class, the trained pool produces Boolean activations
$z_{c,j}(x)$ and signed class sums
\begin{equation}
 f_c(x)=\sum_{j=1}^{K} w_{c,j}z_{c,j}(x),\qquad
 \hat y_{\mathrm{TM}}(x)=\arg\max_c f_c(x).
 \label{eq:tm}
\end{equation}
Flattening class and clause indices gives the activation matrix
$Z\in\{0,1\}^{n\times p}$, initially with $p=CK$.

\begin{figure*}[!t]
\centering
\includegraphics[width=1\textwidth]{fig_cpss_protocol.png}
\caption{CAS pipeline with optional pruning, CPSS, and signed attribution.
Disjoint fit, selection, and report partitions separate signature
extraction, operating-point selection, and evaluation.}
\label{fig:cpss}
\end{figure*}

\subsection{Clause-Pool Pruning}
\label{sec:prune}
Let $m_{c,j}\in\{0,1\}$ retain a clause and
$f_c(x;m)=\sum_j m_{c,j}w_{c,j}z_{c,j}(x)$. The target is a small
pool whose macro-F1 on the selection partition remains within
$\varepsilon$ of the full model:
\begin{equation}
 \min_m\|m\|_0\quad\text{s.t.}\quad
 F_{1,\mathrm{sel}}(m)\ge F_{1,\mathrm{sel}}(\mathbf1)-\varepsilon.
 \label{eq:prune}
\end{equation}
We use a magnitude surrogate rather than solve this discrete objective.
For removed indices $D_c$, set
$R_c=\sum_{j\in D_c}|w_{c,j}|$ and $R=\max_cR_c$.
Since activations are Boolean,
\begin{equation}
 |f_c(x)-f_c(x;m)|\le R_c\le R.
 \label{eq:perturb}
\end{equation}
If $a=\hat y_{\mathrm{TM}}(x)$ and
$\gamma(x)=f_a(x)-\max_{c\ne a}f_c(x)$, the margin of $a$ over any
competing class decreases by at most $2R$. Thus
\begin{equation}
 \gamma(x)>2R\quad\Longrightarrow\quad
 \arg\max_c f_c(x;m)=\hat y_{\mathrm{TM}}(x).
 \label{eq:margin}
\end{equation}
For a fixed deletion budget $k$ per class, the smallest-magnitude
weights minimize this surrogate:
\begin{equation}
 D_c^*(k)=\operatorname{arg\,min}_{|D|=k}
                   \sum_{j\in D}|w_{c,j}|.
 \label{eq:greedy}
\end{equation}


\subsection{CPSS and Signed Evidence}
\label{sec:cpss}
For each class $\ell$, CPSS divides the fit partition into $B$
complementary pairs. On each of the $2B$ halves, a class-balanced
$\ell_1$ logistic model predicts the one-versus-rest class label from
$Z$. For half $h$ and inverse-regularisation strength
$\kappa\in\Lambda$, let $\beta_{\ell j}^{h,\kappa}$ be the fitted
coefficient. A clause receives one vote if any grid point selects it,
using numerical tolerance $\delta$:
\begin{equation}
 \hat\pi_{\ell j}=\frac{1}{2B}\sum_{h=1}^{2B}
 \mathbf1\!\left\{\max_{\kappa\in\Lambda}
             |\beta_{\ell j}^{h,\kappa}|>\delta\right\}.
 \label{eq:stability}
\end{equation}
Let $\bar\beta_{\ell j}$ average coefficients over all halves and
grid points, including zeros. The signed supports are
\begin{equation}
 S_\ell^\pm=\{j:\hat\pi_{\ell j}\ge\pi_{\mathrm{thr}},\quad
                                  \pm\bar\beta_{\ell j}>0\}.
 \label{eq:sign}
\end{equation}
Positive coefficients provide inculpatory evidence; negative ones
provide exculpatory evidence. The attribution rule is
\begin{equation}
 s_\ell(x)=\frac{\sum_{j\in S_\ell^+}\hat\pi_{\ell j}z_j(x)
                 -\sum_{j\in S_\ell^-}\hat\pi_{\ell j}z_j(x)}
                {\sum_{j\in S_\ell^+}\hat\pi_{\ell j}}.
 \label{eq:score}
\end{equation}
Prediction uses $\hat y(x)=\arg\max_\ell s_\ell(x)$.

\subsection{Scope of Stability Error Control}
\label{sec:cert}
For a fixed pool and selector, let $p_j$ denote a clause's population
selection probability on a half-sample, $L_\theta=\{j:p_j\le\theta\}$,
and $q$ the expected number selected per half. Under the sampling
conditions of CPSS~\cite{shah2013}, the sign-agnostic support $S_\ell$
satisfies, for $\pi_{\mathrm{thr}}>1/2$,
\begin{equation}
 \mathbb E|S_\ell\cap L_\theta|
 \le\frac{\theta q}{2\pi_{\mathrm{thr}}-1}.
 \label{eq:bound}
\end{equation}
Taking $\theta=q/p$ yields $q^2/((2\pi_{\mathrm{thr}}-1)p)$.
This controls low-selection-probability inclusions. Interpreting them
as false discoveries requires additional assumptions; using
$\hat q=\sum_j\hat\pi_{\ell j}$ is a plug-in diagnostic, not a
validated false-discovery certificate. Correlated flows, adaptive
operating-point selection, and coefficient signs require separate
consideration. No finite bound is claimed at $\pi_{\mathrm{thr}}=0.5$.

\section{Experimental Setup and Results}
\label{sec:setup}
CICIoMT2024 contains 18 attacks on 40 devices, including 25 real
and 15 simulated devices~\cite{dadkhah2024}. We use its WiFi/MQTT
flow features and six labels: Benign, DDoS, DoS, MQTT, Recon, and
Spoofing. No IP or port fields are included in the selected features.
Sampling is capped at 40,000 training and 15,000 test rows per class.
The training set contains 195,338 rows (19,291 Benign, 16,047
Spoofing, 40,000 in each remaining class); the test set contains
80,868 rows (5,868 Spoofing and 15,000 in each other class).
The 22 numeric features use up to ten quantile bins, giving a cached
154-bit representation that includes redundant bits from repeated
quantiles. The TM uses $K=400$, $T=320$, $s=8$, and 25 epochs.
CPSS uses class-balanced liblinear fits with 200 maximum iterations,
$\delta=10^{-8}$, and $\Lambda=\operatorname{geomspace}(0.001,0.1,6)$.

A fixed permutation (NumPy RandomState seed 0) creates 40,434
CPSS-fit rows, 20,217 selection rows, and 20,217 report rows.
Pruning sweeps $k\in\{0,5,\ldots,395\}$ with $\varepsilon=0.01$.
For each pool, $(B,\pi_{\mathrm{thr}})$ maximizes mean selection
macro-F1 over $B\in\{5,10,15,20,25\}$ and
$\pi_{\mathrm{thr}}\in\{0.5,0.6,0.7,0.8,0.9\}$.


\subsection{Closed-Set Attribution and Sign Validation}
\label{sec:taska}\label{sec:validation}
Full-pool CAS improves over TM argmax for all
three seeds, and the obtained macro-F1 is reported in Table~\ref{tab:headline}. At the same full-pool settings ($B=15$, $\pi_{\mathrm{thr}}=0.8$),
a sign-agnostic scorer that normalizes the stability-weighted sum
over all selected clauses instead yields $0.025\pm0.032$. With full-pool CAS,
mean Benign F1 rises from 0.390 to 0.750, MQTT from 0.703 to 0.824,
and Spoofing from 0.537 to 0.621. DDoS decreases from 0.822 to 0.779
and DoS from 0.833 to 0.751, while Recon remains near 0.869. Thus the macro-F1 improvement reflects substantial
gains for previously weaker classes, rather than uniform gains across
all families.

\begin{table}[t]
\centering
\caption{Comparison for state-of-the-art TM to proposed CAS method. macro-F1 reported as mean$\pm$Standard deviation.}
\label{tab:headline}
\footnotesize
\setlength{\tabcolsep}{3pt}
\begin{tabular}{@{}lccc@{}}
\toprule
Method & $(B,\pi_{\mathrm{thr}})$ & macro-F1 & Accuracy \\
\midrule
TM, full pool & --- & $0.692\pm0.004$ & $0.724\pm0.006$ \\
TM, pruned pool & --- & $0.688\pm0.007$ & $0.720\pm0.010$ \\
CAS, full pool & $(15,0.8)$ & $0.765\pm0.029$ & $0.782\pm0.032$ \\
CAS, pruned pool & $(10,0.9)$ & $\mathbf{0.787\pm0.023}$ & $\mathbf{0.801\pm0.024}$ \\
\bottomrule
\end{tabular}
\end{table}

\subsection{Sensitivity Before and After Pruning}
\label{sec:sweep}
At $B=15$, the pruned-pool threshold sweep spans macro-F1
0.777--0.779 and accuracy 0.792--0.794 (Fig.~\ref{fig:sweep}).
At $\pi_{\mathrm{thr}}=0.8$, varying $B$ from 5 to 25 spans
macro-F1 0.776--0.778 and accuracy 0.790--0.793.
These are fixed sensitivity slices; the selected pruned operating
point remains $(10,0.9)$. The separate full-pool $\kappa$ sweep
(Table~\ref{tab:grid}) gives 0.724--0.763 over the first five grid
points. At threshold 0.8, union-over-grid macro-F1 is 0.765 versus
at most 0.752 for a single grid point.

\begin{figure}[!htb]
\centering
\includegraphics[width=\columnwidth]{fig_sweep_full_pruned_column.png}
\caption{Full (dashed) and pruned (solid) CAS, report-quarter mean$\pm$SD over three seeds:
(a) threshold sweep at $B=15$; (b) pair-count sweep at
$\pi_{\mathrm{thr}}=0.8$. Dotted lines show the pruned-TM baselines.}
\label{fig:sweep}
\end{figure}

\begin{table}[!htb]
\centering
\caption{Parameter grid search for CAS + full pool configuration.}
\label{tab:grid}
\footnotesize
\begin{tabular}{@{}lccccc@{}}
\toprule
$\kappa\backslash\pi_{\mathrm{thr}}$ & 0.5 & 0.6 & 0.7 & 0.8 & 0.9 \\
\midrule
$1.0\!\times\!10^{-3}$ & 0.728 & 0.725 & 0.724 & 0.724 & 0.727 \\
$2.5\!\times\!10^{-3}$ & 0.748 & 0.745 & 0.739 & 0.741 & 0.746 \\
$6.3\!\times\!10^{-3}$ & 0.751 & 0.746 & 0.737 & 0.740 & 0.763 \\
$1.6\!\times\!10^{-2}$ & 0.748 & 0.742 & 0.752 & 0.752 & 0.745 \\
$4.0\!\times\!10^{-2}$ & 0.743 & 0.752 & 0.747 & 0.742 & 0.755 \\
$1.0\!\times\!10^{-1}$ & 0.673 & 0.693 & 0.718 & 0.722 & 0.694 \\
\midrule
Union grid & 0.708 & 0.687 & 0.743 & $\mathbf{0.765}$ & 0.762 \\
\bottomrule
\end{tabular}
\end{table}
\subsection{Pruning and Signature Size}
\label{sec:pruned}\label{sec:signatures}
Figure~\ref{fig:pruning} shows how the weighted TM's report-quarter
macro-F1 and accuracy change as clauses are removed. The curves remain
relatively flat over a broad pruning range, then deteriorate under
more aggressive removal. The selection-quarter rule chooses
$k^*=335$, leaving 65 clauses per class and 390 clauses in total,
a reduction of 83.75\%. This budget is chosen independently of the
report-quarter curves. At the retained budget, the pruned TM achieves
macro-F1 $0.688\pm0.007$ and accuracy $0.720\pm0.010$
(Table~\ref{tab:headline}).

CAS is then applied to the retained 390-clause pool. At the
selection-chosen setting $(B,\pi_{\mathrm{thr}})=(10,0.9)$,
report-quarter macro-F1 reaches $0.787\pm0.023$ and accuracy reaches
$0.801\pm0.024$ (Table~\ref{tab:headline}). From the Fig.~\ref{fig:sweep}, we can understand that removing low-weight clauses may reduce
redundant candidates during signature selection, although these
experiments do not establish the cause of the observed improvement. 

\begin{figure}[!htb]
\centering
\includegraphics[width=0.70\columnwidth]{clause_pruning.png}
\caption{Report-quarter pruning sensitivity (three-seed mean$\pm$SD).
The budget $k^*=335$ is chosen on the separate selection quarter.}
\label{fig:pruning}
\end{figure}

Table~\ref{tab:reduction} compares the full and pruned CAS signatures
under identical CPSS settings for seed 42, where $|S|$ denotes the
total selected support and $|S^+|$ denotes its positive-evidence subset.
Pruning reduces the total support for all six classes; the largest
absolute decrease occurs for DoS, from 248 to 83 clauses, followed by
Spoofing, from 188 to 76. Positive supports decrease for five classes,
but the DoS positive support increases from 32 to 42 despite its smaller
total support. This shows that rerunning CPSS on the retained pool
changes the balance of selected evidence rather than simply shrinking
each signature proportionally.

\begin{table}[!htb]
\centering
\caption{Signature sizes of full and pruned proposed CAS method.}
\label{tab:reduction}
\footnotesize
\begin{tabular}{@{}lcccc@{}}
\toprule
& \multicolumn{2}{c}{Full (2,400)} & \multicolumn{2}{c}{Pruned (390)} \\
\cmidrule(lr){2-3}\cmidrule(lr){4-5}
Class & $|S|$ & $|S^+|$ & $|S|$ & $|S^+|$ \\
\midrule
Benign & 167 & 42 & 63 & 30 \\
DDoS & 75 & 40 & 66 & 36 \\
DoS & 248 & 32 & 83 & 42 \\
MQTT & 145 & 43 & 64 & 26 \\
Recon & 68 & 28 & 58 & 22 \\
Spoofing & 188 & 50 & 76 & 28 \\
\bottomrule
\end{tabular}
\end{table}

\subsection{Special case: Binary Attack Detection}
\begingroup
\color{blue}
\label{sec:binary}
Merging the five attack families gives Benign/Attack labels. The
collapse makes the training prior 90.1\% Attack, so the Benign clause
bank receives Type~II feedback on almost every sample; we oversample
Benign to parity before training. A 188-configuration grid, scored by
per-configuration epoch-checkpoint mean on the selection quarter,
selects 1{,}000 clauses per class (2{,}000 candidates), $T=1000$,
$s=3$, and ten epochs; features and partitions are unchanged.
Selection-quarter pruning with $\varepsilon=0.01$ removes $k^*=130$
clauses per class, retaining 1{,}740 candidates (13.0\% reduction),
and CPSS selects $(B,\pi_{\mathrm{thr}})=(15,0.9)$ and $(5,0.9)$ for
the full and pruned pools. Table~\ref{tab:binary} shows CAS
outperforming its corresponding TM in both metrics, with pruned CAS
slightly exceeding full CAS as in the multiclass case.

\begin{table}[!htb]
\color{blue}
\centering
\caption{Binary classification on the report quarter, mean$\pm$SD
over three seeds. CPSS settings are selected on the selection quarter.}
\label{tab:binary}
\footnotesize
\setlength{\tabcolsep}{3pt}
\begin{tabular}{@{}lccc@{}}
\toprule
Method & $(B,\pi_{\mathrm{thr}})$ & macro-F1 & Accuracy \\
\midrule
TM, full pool & --- & $0.739\pm0.006$ & $0.849\pm0.001$ \\
TM, pruned pool & --- & $0.736\pm0.007$ & $0.848\pm0.002$ \\
CAS, full pool & $(15,0.9)$ & $0.824\pm0.002$ & $0.874\pm0.002$ \\
CAS, pruned pool & $(5,0.9)$ & $\mathbf{0.828\pm0.001}$ & $\mathbf{0.876\pm0.003}$ \\
\bottomrule
\end{tabular}
\end{table}

Table~\ref{tab:binary-reduction} follows Table~\ref{tab:reduction}
at matched settings for seed 42: pruning reduces the Benign support
from 74 to 69 and the Attack support from 77 to 72, both under 80 of
2{,}000 candidates. These counts differ from the selected operating
points in Table~\ref{tab:binary}.

\begin{table}[!htb]
\color{blue}
\centering
\caption{Binary signature sizes for seed 42 at matched settings
$B=15$, $\pi_{\mathrm{thr}}=0.8$.}
\label{tab:binary-reduction}
\footnotesize
\begin{tabular}{@{}lcccc@{}}
\toprule
& \multicolumn{2}{c}{Full (2{,}000)} & \multicolumn{2}{c}{Pruned (1{,}740)} \\
\cmidrule(lr){2-3}\cmidrule(lr){4-5}
Class & $|S|$ & $|S^+|$ & $|S|$ & $|S^+|$ \\
\midrule
Benign & 74 & 51 & 69 & 45 \\
Attack & 77 & 25 & 72 & 24 \\
\bottomrule
\end{tabular}
\end{table}

Full CAS raises attack precision from 0.894 to 0.974 and Benign recall
from 0.529 to 0.902, lowers attack recall from 0.923 to 0.868, and
raises attack F1 from 0.908 to 0.918. Its signed attack score reaches
AUROC $0.929\pm0.021$, against $0.874\pm0.029$ for the TM class-sum
margin, 0.853 unclipped and 0.827 clipped to $[0,T]$; pruned CAS
matches it at $0.929\pm0.023$.

Figure~\ref{fig:binary-sweep} follows Fig.~\ref{fig:sweep}: the
threshold and pair-count sweeps give full/pruned macro-F1 ranges
0.814--0.824/0.816--0.826 and 0.819--0.823/0.825--0.827, the pruned
pool being at least as stable in both slices.

\begin{figure}[!htb]
\color{blue}
\centering
\includegraphics[width=0.84\columnwidth]{fig_binary_sweep_full_pruned_column.pdf}
\caption{Binary CAS sensitivity, matching Fig.~\ref{fig:sweep}:
full (dashed) and pruned (solid) report-quarter mean$\pm$SD over three
seeds. (a) Threshold sweep at $B=15$; (b) pair-count sweep at
$\pi_{\mathrm{thr}}=0.8$. Dotted lines show pruned-TM baselines.}
\label{fig:binary-sweep}
\end{figure}

Figure~\ref{fig:binary-pruning} follows Fig.~\ref{fig:pruning}. At
$k^*=130$ the budget is near-free: TM macro-F1 moves from 0.739 to
0.736, accuracy from 0.849 to 0.848, and Benign/Attack F1 from
0.569/0.908 to 0.564/0.908. Beyond roughly 700 removed clauses per
class both degrade sharply.

\begin{figure}[!htb]
\color{blue}
\centering
\includegraphics[width=0.58\columnwidth]{fig_binary_clause_pruning.pdf}
\caption{Binary TM pruning: macro-F1, accuracy, and Benign/Attack F1
(report-quarter mean$\pm$SD, three seeds). The selection-quarter
budget $k^*=130$ retains 870 clauses per class (1{,}740 total).}
\label{fig:binary-pruning}
\end{figure}
\section{Discussion and Conclusion}
CAS provides signed, inspectable evidence. Multiclass pruning improves
CAS macro-F1 from 0.765 to 0.787 and reduces threshold sensitivity
(Fig.~\ref{fig:sweep}). Binary pruning likewise raises CAS macro-F1 from
0.824 to 0.828 while preserving TM macro-F1
(Figs.~\ref{fig:binary-sweep}--\ref{fig:binary-pruning}), so the
benefit of compactness holds in both tasks once the binary detector is
trained on a balanced prior. Full binary CAS improves attack precision
from 0.894 to 0.974 and Benign recall from 0.529 to 0.902, but raises
missed attacks from 7.7\% to 13.2\%, requiring cost-aware
operating-point selection.

The margin condition and CPSS bound do not certify causal explanations.
This retrospective study covers three seeds and one capped dataset,
with prior test inspection and no matched XAI baseline. Independent
devices and time periods, semantic rule stability, analyst usefulness,
and runtime and energy measurements remain necessary for deployment
assessment.

\endgroup

% All technical content and floats must finish by page 4.
% A subsequent page is reserved exclusively for references.
\clearpage

\begin{thebibliography}{12}
\bibitem{yin2017} C. Yin, Y. Zhu, J. Fei, and X. He,
``A Deep Learning Approach for Intrusion Detection Using Recurrent
Neural Networks,'' \emph{IEEE Access}, vol. 5, pp. 21954--21961, 2017.
\bibitem{xaiids2025} P. Hermosilla, S. Berr\'ios, and H. Allende-Cid,
``Explainable AI for Forensic Analysis: A Comparative Study of SHAP and
LIME in Intrusion Detection Models,'' \emph{Applied Sciences},
vol. 15, no. 13, art. 7329, 2025.
\bibitem{granmo2018} O.-C. Granmo, ``The Tsetlin Machine --- a
game-theoretic bandit-driven approach to optimal pattern recognition with
propositional logic,'' \emph{arXiv:1804.01508}, 2018.
\bibitem{phoulady2020} A. Phoulady, O.-C. Granmo, S. R. Gorji, and H. A. Phoulady,
``The Weighted Tsetlin Machine: Compressed Representations with Weighted
Clauses,'' \emph{arXiv:1911.12607}, 2020.
\bibitem{abeyrathna2020} K. D. Abeyrathna et al., ``Intrusion Detection
with Interpretable Rules Generated using the Tsetlin Machine,'' \emph{IEEE
SSCI}, 2020.
\bibitem{jaiswal2026iomt} R. Jaiswal, P.-A. Andersen, L. R. Cenkeramaddi,
L. Jiao, and O.-C. Granmo, ``A Tsetlin Machine-driven Intrusion Detection
System for Next-Generation IoMT Security,'' \emph{IEEE SVCC} /
\emph{arXiv:2604.03205}, 2026.
\bibitem{jaiswal2026medsec} R. Jaiswal et al., ``On-Device Interpretable
Tsetlin Machine-Based Intrusion Detection for Secure IoMT,''
\emph{arXiv:2605.16707}, 2026.
\bibitem{meinshausen2010} N. Meinshausen and P. B\"uhlmann, ``Stability
selection,'' \emph{JRSS-B}, vol. 72, no. 4, pp. 417--473, 2010.
\bibitem{shah2013} R. D. Shah and R. J. Samworth, ``Variable selection with error
control: another look at stability selection,'' \emph{JRSS-B}, vol. 75,
no. 1, pp. 55--80, 2013.
\bibitem{tmu} CAIR, ``Tsetlin Machine Unified (TMU),'' \url{https://github.com/cair/tmu}.
\bibitem{dadkhah2024} S. Dadkhah et al., ``CICIoMT2024: A benchmark dataset
for multi-protocol security assessment in IoMT,'' \emph{Internet of Things},
vol. 28, art. 101351, 2024.
\bibitem{du1993} D.-Z. Du and F. K. Hwang, \emph{Combinatorial Group Testing
and Its Applications}, World Scientific, 1993.
\end{thebibliography}

\end{document}
